% ─────────────────────────────────────────────────────────────────────────────
%
% Especificación y decisiones de diseño del módulo spatial_awareness.
% Documento base para la implementación y para el informe de tesis.
% ─────────────────────────────────────────────────────────────────────────────

\section{Conciencia espacial de obstáculos fuera del campo comunicado}
\label{sec:spatial-awareness}

El sistema de percepción desarrollado hasta esta etapa comunica al usuario los
obstáculos contenidos en una grilla frontal generada a partir del fotograma de
profundidad actual. La apertura horizontal de esa grilla puede configurarse en
tiempo de ejecución. Como consecuencia, existen obstáculos que fueron
observados e incorporados al mapa local, pero que dejan de estar representados
en la señal háptica frontal cuando quedan a los lados o detrás del usuario.

Esta situación es especialmente relevante durante el desplazamiento. Por
ejemplo, una pared inicialmente observada puede quedar fuera del campo
comunicado cuando el usuario gira la cabeza y, aun así, continuar
aproximándose lateralmente a ella. El flujo frontal normal deja de advertirla,
aunque el obstáculo permanece disponible en el mapa de ocupación acumulado.

Para cubrir este caso se propone el paquete \texttt{spatial\_awareness}. Su
objetivo es detectar riesgo de colisión con celdas ocupadas que se encuentran
fuera del campo angular ya comunicado al usuario. El nodo no sustituye la
representación frontal ni vuelve a señalizar obstáculos contenidos en ella.
Opera únicamente como una capa complementaria para las regiones izquierda,
derecha y trasera.

La activación de una advertencia requiere simultáneamente que:
%
\begin{enumerate}
  \item exista una celda o agrupación de celdas ocupadas fuera del campo
    comunicado;
  \item la distancia entre la cámara y el obstáculo esté dentro del intervalo
    operativo;
  \item la velocidad del usuario tenga una componente positiva en dirección
    al obstáculo; y
  \item el mapa y la odometría sean recientes y válidos.
\end{enumerate}

La intensidad se determina principalmente mediante el tiempo estimado hasta
la colisión (\emph{time to collision}, TTC). De esta manera, distancia y
velocidad no se tratan como factores independientes arbitrarios: ambas
variables se combinan en una magnitud cinemática directamente relacionada con
la urgencia de la advertencia.

\subsection{Alcance y restricciones}
\label{subsec:sa-scope}

La primera implementación queda limitada a las siguientes condiciones:
%
\begin{itemize}
  \item la fuente de odometría es RTAB-Map;
  \item el cálculo horizontal utiliza la convención interna del proyecto:
    $X$ hacia la derecha, $Y$ hacia abajo y $Z$ hacia adelante;
  \item el mapa de entrada es el
    \texttt{/local\_mapper/occupancy\_grid}, expresado en el marco
    \texttt{odom};
  \item el campo comunicado es horizontal, simétrico respecto del eje frontal
    de la cámara;
  \item el módulo considera únicamente celdas ya presentes en el mapa local;
    no puede detectar obstáculos que nunca hayan sido observados por la
    cámara; y
  \item la antigüedad individual de cada celda no se evalúa en esta etapa. La
    gestión temporal de las observaciones se incorporará posteriormente en el
    \texttt{local\_mapper}, sin modificar la responsabilidad de
    \texttt{spatial\_awareness}.
\end{itemize}

El nombre del paquete y del nodo es \texttt{spatial\_awareness}. Se descarta
la variante \texttt{spatial\_awarness}, utilizada inicialmente en el nombre de
la rama, por contener un error ortográfico.

\subsection{Arquitectura}
\label{subsec:sa-architecture}

El módulo se implementará como un nodo ROS~2 independiente en C++17. La
separación respecto de \texttt{local\_mapper} permite mantener responsabilidades
claras: el mapeador estima ocupación y \texttt{spatial\_awareness} interpreta
ese mapa en función del movimiento del usuario y del campo comunicado.

El flujo de datos propuesto es:
%
\begin{center}
\texttt{RTAB-Map}
$\longrightarrow$
\texttt{nav\_odom}
$\longrightarrow$
\texttt{spatial\_awareness}

\texttt{local\_mapper}
$\longrightarrow$
\texttt{occupancy\_grid}
$\longrightarrow$
\texttt{spatial\_awareness}

\texttt{aperture\_deg}
$\longrightarrow$
\texttt{spatial\_awareness}
$\longrightarrow$
\texttt{collision\_risk}
\end{center}

La Tabla~\ref{tab:sa-topics} resume las interfaces previstas.

\begin{table}[ht]
  \centering
  \caption{Topics del nodo \texttt{spatial\_awareness}.}
  \label{tab:sa-topics}
  \small
  \begin{tabular}{llp{7cm}}
    \hline
    \textbf{Dirección} & \textbf{Topic y tipo} & \textbf{Descripción} \\
    \hline
    Sub &
    \texttt{/local\_mapper/occupancy\_grid}
    (\texttt{nav\_msgs/OccupancyGrid}) &
    Mapa local acumulado en el marco \texttt{odom}. \\

    Sub &
    \texttt{nav\_odom}
    (\texttt{nav\_msgs/Odometry}) &
    Pose y velocidad provenientes de RTAB-Map. \\

    Sub &
    \texttt{/perception/depth\_grid/aperture\_deg}
    (\texttt{std\_msgs/Float32}) &
    Semiángulo horizontal comunicado al usuario. \\

    Pub &
    \texttt{/spatial\_awareness/collision\_risk}
    (\texttt{custom\_interfaces/SpatialAwareness}) &
    Riesgo independiente hacia izquierda, derecha y atrás. \\
    \hline
  \end{tabular}
\end{table}

El estado de apertura debe estar disponible también para nodos que se inicien
después de haber enviado el último comando. Por ello, se recomienda publicar
el valor efectivo con QoS \texttt{transient\_local}, o introducir un topic de
estado separado del topic de comando. El valor consumido por
\texttt{spatial\_awareness} debe ser el valor aplicado por el encoder, no
solamente la última solicitud emitida por una interfaz gráfica.

\subsubsection{Interfaz de salida}
\label{subsubsec:sa-output}

El consumidor futuro necesita distinguir izquierda, derecha y atrás, junto
con la intensidad correspondiente. Se propone un único mensaje con tres
canales para garantizar que los estados pertenezcan al mismo instante y evitar
la sincronización de tres topics independientes.

La interfaz se divide en dos mensajes:

\begin{verbatim}
# custom_interfaces/msg/DirectionalRisk.msg
bool active
float32 intensity
float32 distance_m
float32 closing_speed_mps
float32 time_to_collision_s
\end{verbatim}

\begin{verbatim}
# custom_interfaces/msg/SpatialAwareness.msg
std_msgs/Header header
bool valid
DirectionalRisk left
DirectionalRisk right
DirectionalRisk rear
\end{verbatim}

El campo \texttt{intensity} se expresa en el intervalo $[0,100]$. Los campos
de distancia, velocidad de cierre y TTC no son necesarios para el actuador
final, pero se conservan para diagnóstico, visualización, calibración y
reproducibilidad experimental.

\texttt{valid=false} indica que no puede calcularse un estado confiable por
falta o antigüedad excesiva del mapa o de la odometría. Cuando los datos son
válidos pero no existe riesgo, \texttt{valid=true} y los tres canales se
publican con \texttt{active=false} e intensidad nula.

\subsection{Convenciones de coordenadas}
\label{subsec:sa-frames}

RTAB-Map utiliza externamente la convención REP-103: $X$ hacia adelante, $Y$
hacia la izquierda y $Z$ hacia arriba. El pipeline de percepción del proyecto
trabaja internamente con la convención óptica:
%
\begin{equation}
  X_{\text{int}} \text{ hacia la derecha}, \qquad
  Y_{\text{int}} \text{ hacia abajo}, \qquad
  Z_{\text{int}} \text{ hacia adelante}.
\end{equation}

La conversión ya utilizada por \texttt{depth\_obstacle\_filter} para la
posición es:
%
\begin{equation}
  \label{eq:sa-rtab-position}
  \mathbf{p}_{\text{int}}
  =
  \begin{bmatrix}
    -p_{\text{ros},y} \\
    -p_{\text{ros},z} \\
     p_{\text{ros},x}
  \end{bmatrix}.
\end{equation}

\texttt{spatial\_awareness} debe aplicar la misma conversión antes de realizar
cualquier cálculo geométrico. De esta forma, una velocidad positiva en $X$
representa desplazamiento hacia la derecha y una velocidad positiva en $Z$
representa desplazamiento hacia adelante.

Para clasificar obstáculos respecto de la cabeza se necesita además el yaw de
la cámara. Sea $\mathbf{q}$ la orientación convertida al convenio interno. El
vector frontal en el marco \texttt{odom} se obtiene como:
%
\begin{equation}
  \label{eq:sa-forward}
  \mathbf{f}_{xz}
  =
  \operatorname{normalize}
  \left(
    \Pi_{xz}
    \left[
      \mathbf{q}
      \begin{pmatrix}0\\0\\1\end{pmatrix}
      \mathbf{q}^{*}
    \right]
  \right),
\end{equation}
%
donde $\Pi_{xz}$ proyecta el vector sobre el plano horizontal XZ. El vector
hacia la derecha se define como:
%
\begin{equation}
  \label{eq:sa-right}
  \mathbf{r}_{xz} =
  \begin{pmatrix}f_z\\-f_x\end{pmatrix}.
\end{equation}

\subsection{Estimación de velocidad}
\label{subsec:sa-velocity}

El campo \texttt{twist} de \texttt{nav\_msgs/Odometry} es la fuente
convencional de velocidad para un consumidor ROS~2. Por lo tanto, la primera
opción será utilizar la velocidad lineal publicada por RTAB-Map, siempre que:
%
\begin{itemize}
  \item sus componentes sean finitas;
  \item el mensaje tenga una antigüedad menor al timeout configurado;
  \item la norma se encuentre dentro de límites físicamente plausibles para
    una persona caminando; y
  \item se haya verificado experimentalmente el frame en el que RTAB-Map
    expresa el \texttt{twist}.
\end{itemize}

Si la velocidad publicada no está disponible, es nula de forma persistente
durante desplazamiento, presenta discontinuidades o su frame no puede
establecerse de forma confiable, se utilizará como respaldo una estimación a
partir de diferencias de posición:
%
\begin{equation}
  \label{eq:sa-velocity-raw}
  \widehat{\mathbf{v}}_t =
  \frac{\mathbf{p}_t-\mathbf{p}_{t-1}}{t-t_{t-1}}.
\end{equation}

La velocidad se calcula en el marco \texttt{odom}, no en el marco instantáneo
de la cámara. Esto permite compararla directamente con la dirección desde el
usuario hacia cada celda del mapa.

Para reducir el ruido de la diferenciación se aplica un filtro pasa bajos
exponencial:
%
\begin{equation}
  \label{eq:sa-velocity-filter}
  \mathbf{v}_t =
  \alpha_v\,\widehat{\mathbf{v}}_t
  +(1-\alpha_v)\,\mathbf{v}_{t-1},
  \qquad 0 < \alpha_v \leq 1.
\end{equation}

El parámetro \texttt{min\_motion\_speed\_mps} define la velocidad mínima para
considerar que el usuario está desplazándose. Se propone inicialmente
$0{,}10\,\text{m/s}$: es suficientemente menor que una velocidad normal de
marcha, típicamente cercana a $1\,\text{m/s}$, pero evita activar
advertencias por jitter de odometría cuando la persona permanece quieta. El
valor quedará configurable y deberá calibrarse con registros reales.

\subsection{Selección de celdas ocupadas}
\label{subsec:sa-occupied-cells}

Cada elemento conocido del \texttt{OccupancyGrid} tiene un valor entre 0 y
100. Una celda se considera candidata cuando:
%
\begin{equation}
  \label{eq:sa-occupancy-threshold}
  m_i \geq P_{\text{occ,min}} = 65.
\end{equation}

Las celdas desconocidas, con valor $-1$, se descartan. El umbral inicial de 65
se mantiene como parámetro porque deberá calibrarse junto con el modelo
log-odds del mapa.

Para rechazar impactos aislados se exige que cada candidato pertenezca a un
componente conexo con un número mínimo de celdas ocupadas. Se utilizará
conectividad de ocho vecinos para no separar superficies diagonales. El
parámetro propuesto es:
%
\begin{equation}
  N_{\text{cluster,min}} = 3.
\end{equation}

Con una resolución de $0{,}10\,\text{m}$, tres celdas representan evidencia
espacial suficiente para rechazar un punto aislado sin exigir una superficie
grande. Tanto el valor como el tipo de conectividad quedarán encapsulados en
la librería de cálculo para permitir su validación unitaria.

La posición de una celda se calcula en su centro geométrico y se transforma
mediante la pose completa de \texttt{info.origin}. No debe asumirse que el
\texttt{OccupancyGrid} está contenido en el plano XY sin rotación: el mapa
actual aplica una rotación de $+90^\circ$ alrededor de $X$ para representar
el plano XZ.

\subsection{Campo comunicado}
\label{subsec:sa-fov}

El parámetro \texttt{h\_aperture\_deg} conserva la semántica actual de
semiángulo. Por ejemplo:
%
\begin{equation}
  h_{\text{aperture}}=45^\circ
  \quad\Longrightarrow\quad
  \text{campo comunicado}=[-45^\circ,+45^\circ],
\end{equation}
%
es decir, una apertura horizontal total de $90^\circ$.

Para una celda ubicada en $\mathbf{o}_i$ y una cámara ubicada en
$\mathbf{p}$, se define el vector horizontal relativo:
%
\begin{equation}
  \label{eq:sa-relative}
  \mathbf{d}_i =
  \begin{pmatrix}
    o_{i,x}-p_x\\
    o_{i,z}-p_z
  \end{pmatrix}.
\end{equation}

Su ángulo firmado respecto de la dirección frontal es:
%
\begin{equation}
  \label{eq:sa-angle}
  \theta_i =
  \operatorname{atan2}
  \left(
    \mathbf{r}_{xz}\cdot\mathbf{d}_i,\,
    \mathbf{f}_{xz}\cdot\mathbf{d}_i
  \right).
\end{equation}

Una celda pertenece al campo ya comunicado cuando:
%
\begin{equation}
  \label{eq:sa-in-fov}
  |\theta_i| \leq h_{\text{aperture}} + h_{\text{margin}}.
\end{equation}

El margen angular evita que una misma superficie alterne rápidamente entre el
flujo frontal y \texttt{spatial\_awareness} por ruido de yaw o discretización
de celdas. Se propone inicialmente
$h_{\text{margin}}=3^\circ$. Toda celda que satisface la
ecuación~\ref{eq:sa-in-fov} se descarta del módulo, porque será comunicada por
el flujo frontal normal.

La velocidad angular no modifica directamente la intensidad. Si el usuario
gira y una celda entra en el campo frontal, deja de ser responsabilidad de
\texttt{spatial\_awareness}; si sale del campo y continúa existiendo riesgo
cinemático, pasa a ser considerada por este módulo.

\subsection{Modelo geométrico de colisión}
\label{subsec:sa-collision-model}

El intervalo de distancia operativo se evalúa desde la cámara hasta el centro
de la celda:
%
\begin{equation}
  \label{eq:sa-distance-range}
  d_{\min}=0{,}20\,\text{m}
  \leq d_i = \|\mathbf{d}_i\|
  \leq d_{\max}=1{,}00\,\text{m}.
\end{equation}

Las celdas a más de $1\,\text{m}$ no generan señal. Las celdas a menos de
$0{,}20\,\text{m}$ también se descartan: se encuentran fuera del rango útil
definido para esta advertencia y se considera que la colisión ya ocurrió o
que la medición puede corresponder a ruido muy próximo.

El radio corporal se modela mediante un parámetro
\texttt{body\_radius\_m}, inicialmente igual a $0{,}20\,\text{m}$. Para no
aplicar dos veces la compensación de proximidad, este radio no se resta de
$d_i$ ni modifica los límites de la ecuación~\ref{eq:sa-distance-range}. Se
utiliza para definir un corredor de trayectoria: una celda solo se considera
riesgo de colisión si su separación lateral respecto de la dirección de
movimiento es menor que el radio corporal más una tolerancia de celda.

Sea:
%
\begin{equation}
  \mathbf{\hat v} =
  \frac{\mathbf{v}}{\|\mathbf{v}\|}.
\end{equation}

La proyección longitudinal y la distancia lateral son:
%
\begin{equation}
  \label{eq:sa-corridor}
  s_i = \mathbf{d}_i\cdot\mathbf{\hat v},
  \qquad
  \ell_i =
  \left\|
    \mathbf{d}_i-s_i\mathbf{\hat v}
  \right\|.
\end{equation}

La celda pertenece al corredor de colisión cuando:
%
\begin{equation}
  \label{eq:sa-corridor-condition}
  s_i > 0
  \quad\land\quad
  \ell_i \leq r_{\text{body}} + r_{\text{cell}},
\end{equation}
%
donde $r_{\text{cell}}$ representa la mitad de la diagonal de la celda. Esta
condición evita advertir por obstáculos cercanos hacia los cuales el usuario
no se dirige.

\subsection{Velocidad de cierre y tiempo hasta colisión}
\label{subsec:sa-ttc}

La dirección unitaria desde la cámara hacia la celda es:
%
\begin{equation}
  \label{eq:sa-direction}
  \mathbf{\hat d}_i =
  \frac{\mathbf{d}_i}{\|\mathbf{d}_i\|}.
\end{equation}

La velocidad de cierre se define como la proyección de la velocidad del
usuario sobre esa dirección:
%
\begin{equation}
  \label{eq:sa-closing-speed}
  v_{\text{close},i}
  =
  \max\left(0,\mathbf{v}\cdot\mathbf{\hat d}_i\right).
\end{equation}

Si $v_{\text{close},i}=0$, el usuario se encuentra quieto, se mueve
tangencialmente o se aleja del obstáculo. En los tres casos no se genera
advertencia.

Para evitar activaciones por ruido se exige:
%
\begin{equation}
  \label{eq:sa-closing-threshold}
  v_{\text{close},i} \geq v_{\text{close,min}},
\end{equation}
%
con un valor inicial configurable de $0{,}10\,\text{m/s}$.

El TTC se calcula como:
%
\begin{equation}
  \label{eq:sa-ttc}
  \operatorname{TTC}_i =
  \frac{d_i}{v_{\text{close},i}}.
\end{equation}

Este modelo supone velocidad constante durante el horizonte corto de
advertencia. Para distancias menores a un metro y velocidades peatonales, la
aproximación es suficiente como primer criterio y evita introducir un modelo
de aceleración difícil de estimar de manera robusta con odometría visual.

\subsection{Cálculo de intensidad}
\label{subsec:sa-intensity}

La intensidad debe aumentar cuando disminuye el tiempo disponible para evitar
la colisión. Se propone una relación lineal acotada e inversa respecto del TTC:
%
\begin{equation}
  \label{eq:sa-intensity}
  I_i =
  100\,
  \operatorname{clamp}
  \left(
    \frac{T_{\max}-\operatorname{TTC}_i}
         {T_{\max}-T_{\min}},
    0,\,
    1
  \right).
\end{equation}

$T_{\max}$ representa el horizonte a partir del cual no se emite señal y
$T_{\min}$ el TTC para intensidad máxima. Como valores iniciales se proponen:
%
\begin{equation}
  T_{\max}=2{,}0\,\text{s},
  \qquad
  T_{\min}=0{,}25\,\text{s}.
\end{equation}

La compuerta espacial de $1\,\text{m}$ continúa aplicándose antes de esta
ecuación. Por lo tanto, un obstáculo lejano no genera señal aunque la
estimación de velocidad produzca un TTC bajo por un valor espurio.

Esta formulación conserva el comportamiento requerido:
%
\begin{itemize}
  \item a igual distancia, una mayor velocidad de aproximación produce mayor
    intensidad;
  \item a igual velocidad, una menor distancia produce mayor intensidad;
  \item una velocidad nula o dirigida en sentido opuesto produce intensidad
    nula; y
  \item las celdas fuera del intervalo $[0{,}20,1{,}00]\,\text{m}$ no generan
    señal.
\end{itemize}

Los parámetros $T_{\min}$ y $T_{\max}$ deberán calibrarse mediante rosbag con
trayectorias peatonales reales. Si la respuesta lineal no resulta perceptualmente
adecuada, podrá sustituirse por una curva exponencial o sigmoidea sin cambiar
la arquitectura ni la interfaz.

\subsection{Clasificación direccional}
\label{subsec:sa-directions}

Después de descartar el campo frontal, cada candidato se asigna a una de tres
regiones según el ángulo $\theta_i$:
%
\begin{equation}
  \label{eq:sa-direction-classes}
  \operatorname{region}(\theta_i)=
  \begin{cases}
    \text{derecha} &
      h_{\text{aperture}}+h_{\text{margin}} < \theta_i < 135^\circ,\\
    \text{izquierda} &
      -135^\circ < \theta_i <
      -h_{\text{aperture}}-h_{\text{margin}},\\
    \text{atrás} &
      |\theta_i| \geq 135^\circ.
  \end{cases}
\end{equation}

El límite de $135^\circ$ se propone como valor configurable
$\texttt{rear\_half\_angle\_deg}=45^\circ$, de modo que la región trasera cubra
$90^{\circ}$ centrados en $180^{\circ}$. Esta división entrega una semántica
estable al actuador: izquierda y derecha cubren los laterales, mientras que
atrás representa obstáculos próximos a la dirección opuesta a la mirada.

Para cada región se selecciona la agrupación con mayor intensidad. En caso de
empate se prioriza el menor TTC y luego la menor distancia. El mensaje puede
contener hasta tres advertencias activas simultáneamente, una por región.

\subsection{Estabilidad temporal y fallos}
\label{subsec:sa-robustness}

\subsubsection{Histéresis}

La discretización del mapa y el ruido de velocidad pueden producir
oscilaciones alrededor de los umbrales. Se aplicarán:
%
\begin{itemize}
  \item histéresis de velocidad, con umbral de activación mayor que el de
    desactivación;
  \item margen angular en el borde del campo comunicado;
  \item filtrado pasa bajos de velocidad; y
  \item limitación de la variación máxima de intensidad por segundo.
\end{itemize}

No se mantendrá una señal cuando la velocidad de cierre desaparezca: la
histéresis solo evita conmutaciones por ruido y debe tener una duración corta.

\subsubsection{Datos obsoletos}

Se mantienen timestamps independientes para mapa, odometría y apertura. Si el
mapa o la odometría superan su timeout:
%
\begin{itemize}
  \item no se calculan nuevos riesgos;
  \item se publica, opcionalmente durante un ciclo, un mensaje
    \texttt{valid=false} con todos los canales desactivados; y
  \item no se conserva la última intensidad activa.
\end{itemize}

La recomendación es continuar publicando un estado nulo e inválido a una tasa
baja, en lugar de guardar silencio. Esto permite al consumidor distinguir
entre ausencia de riesgo y ausencia de datos. Si el contrato del futuro
actuador interpreta timeout como apagado seguro, esta política puede
configurarse para dejar de publicar.

\subsubsection{Frecuencia del mapa}

La cámara se utiliza a 6~Hz, por lo que aplicar un \emph{throttle} adicional
al \texttt{OccupancyGrid} introduciría una latencia incompatible con una
advertencia basada en TTC. En consecuencia, el mapa debe publicarse después de
cada actualización, configurando
\texttt{occupancy\_publish\_every\_n\_frames=1}. La tasa efectiva del mapa es
entonces la tasa de la cámara, aproximadamente 6~Hz, con un período de
$167\,\text{ms}$ entre observaciones.

El nodo \texttt{spatial\_awareness} evalúa y publica el riesgo a 10~Hz usando
el mapa y la odometría más recientes. Esta frecuencia no crea información
geométrica nueva entre fotogramas, pero permite reaccionar a actualizaciones
de odometría y emitir con rapidez un estado inválido o nulo ante timeouts. No
se modifica en esta etapa la semántica ni el envejecimiento de las celdas del
mapa.

\subsection{Parámetros configurables}
\label{subsec:sa-parameters}

\begin{table}[ht]
  \centering
  \caption{Parámetros iniciales propuestos para
  \texttt{spatial\_awareness}.}
  \label{tab:sa-parameters}
  \small
  \begin{tabular}{lcp{7cm}}
    \hline
    \textbf{Parámetro} & \textbf{Inicial} & \textbf{Descripción} \\
    \hline
    \texttt{occupancy\_threshold} & 65 &
      Valor mínimo de ocupación de una celda. \\
    \texttt{min\_cluster\_cells} & 3 &
      Celdas ocupadas mínimas por componente conexo. \\
    \texttt{min\_distance\_m} & 0,20 &
      Distancia celda-cámara por debajo de la cual no se señaliza. \\
    \texttt{max\_distance\_m} & 1,00 &
      Alcance máximo de señalización. \\
    \texttt{body\_radius\_m} & 0,20 &
      Radio del corredor corporal; no se resta de la distancia medida. \\
    \texttt{min\_motion\_speed\_mps} & 0,10 &
      Movimiento mínimo para habilitar el análisis. \\
    \texttt{min\_closing\_speed\_mps} & 0,10 &
      Velocidad mínima de aproximación a una celda. \\
    \texttt{closing\_speed\_hysteresis\_mps} & 0,03 &
      Reducción del umbral de desactivación para evitar conmutaciones por ruido. \\
    \texttt{velocity\_filter\_alpha} & 0,25 &
      Coeficiente del filtro exponencial de velocidad. \\
    \texttt{ttc\_max\_s} & 2,00 &
      TTC para intensidad nula. \\
    \texttt{ttc\_min\_s} & 0,25 &
      TTC para intensidad máxima. \\
    \texttt{fov\_margin\_deg} & 3,0 &
      Margen de exclusión alrededor del campo frontal. \\
    \texttt{rear\_half\_angle\_deg} & 45,0 &
      Semiángulo de la región trasera. \\
    \texttt{odom\_timeout\_s} & 0,25 &
      Antigüedad máxima de la odometría. \\
    \texttt{map\_timeout\_s} & 0,30 &
      Antigüedad máxima del mapa para producir una salida válida. \\
    \texttt{intensity\_slew\_rate} & 200/s &
      Variación máxima de intensidad por segundo. \\
    \hline
  \end{tabular}
\end{table}

Todos los valores son hipótesis iniciales de ingeniería y no resultados
experimentales. Deberán distinguirse como tales en el informe final.

\subsection{Descomposición de software}
\label{subsec:sa-software}

Se propone separar el cálculo en una librería sin dependencias directas de
ROS y un nodo de integración:
%
\begin{description}
  \item[\texttt{SpatialRiskEvaluator}] recibe pose 2D, velocidad, apertura y
    una colección de celdas ocupadas. Devuelve los tres riesgos direccionales.
    Contiene la selección por FOV, corredor corporal, velocidad de cierre, TTC
    e intensidad.

  \item[\texttt{OccupancyClusterer}] convierte la grilla en componentes
    conexos y elimina agrupaciones menores que
    \texttt{min\_cluster\_cells}.

  \item[\texttt{VelocityEstimator}] valida el \texttt{twist} de RTAB-Map y
    mantiene el respaldo por diferencia de pose con filtrado temporal.

  \item[\texttt{SpatialAwarenessNode}] administra suscripciones, QoS,
    timestamps, parámetros, conversión de mensajes y publicación.
\end{description}

Esta separación permite probar la geometría y el modelo de riesgo con entradas
deterministas, sin requerir un grafo ROS ni hardware.

\subsection{Plan de implementación}
\label{subsec:sa-plan}

La implementación se divide en las siguientes etapas:
%
\begin{enumerate}
  \item Crear el paquete C++ \texttt{spatial\_awareness} y su configuración.

  \item Agregar \texttt{DirectionalRisk.msg} y
    \texttt{SpatialAwareness.msg} a \texttt{custom\_interfaces}.

  \item Implementar y probar la conversión entre la salida REP-103 de RTAB-Map
    y la convención interna X-derecha, Y-abajo, Z-adelante.

  \item Verificar mediante rosbag la semántica de
    \texttt{nav\_odom.twist}. Implementar el respaldo por diferencia de pose.

  \item Implementar la decodificación general del
    \texttt{OccupancyGrid}, respetando resolución, origen y orientación.

  \item Implementar agrupamiento de celdas, exclusión por FOV y clasificación
    izquierda-derecha-atrás.

  \item Implementar corredor corporal, velocidad de cierre, TTC e intensidad.

  \item Incorporar timeouts, histéresis y limitación de variación.

  \item Integrar el nodo en \texttt{nav\_bringup} con el argumento
    \texttt{use\_spatial\_awareness}.

  \item Publicar marcadores de diagnóstico en RViz para visualizar el campo
    excluido, el corredor de movimiento, las agrupaciones candidatas y el
    riesgo seleccionado en cada dirección.

  \item Configurar el \texttt{local\_mapper} para publicar el mapa en cada
    fotograma de la cámara, aproximadamente a 6~Hz, sin modificar todavía el
    modelo temporal por celda.

  \item Validar primero con escenas sintéticas y luego con rosbag de caminatas
    reales junto a paredes y obstáculos laterales.
\end{enumerate}

\subsection{Estrategia de pruebas}
\label{subsec:sa-tests}

\subsubsection{Pruebas unitarias}

La librería deberá cubrir, como mínimo, los siguientes casos:
%
\begin{enumerate}
  \item pared a la derecha fuera del FOV y movimiento hacia la derecha:
    señal derecha activa;
  \item misma pared con usuario quieto: ninguna señal;
  \item misma pared con movimiento en sentido opuesto: ninguna señal;
  \item obstáculo dentro del FOV: excluido del módulo;
  \item obstáculo a más de $1\,\text{m}$: excluido;
  \item obstáculo a menos de $0{,}20\,\text{m}$: excluido;
  \item celda aislada por debajo del tamaño mínimo de agrupación: excluida;
  \item movimiento diagonal hacia dos regiones: riesgos independientes;
  \item obstáculo trasero y retroceso: señal trasera activa;
  \item giro que hace entrar una pared al FOV: la señal lateral se desactiva;
  \item mapa u odometría obsoletos: mensaje inválido y señales nulas;
  \item igualdad de distancia con distintas velocidades: mayor velocidad
    produce mayor intensidad; y
  \item igualdad de velocidad con distintas distancias: menor distancia
    produce mayor intensidad.
\end{enumerate}

\subsubsection{Pruebas con rosbag}

Se propone registrar los siguientes escenarios:
%
\begin{description}
  \item[E1: pared lateral.] Caminata paralela a una pared con aproximación
    gradual hacia ella, manteniéndola fuera del campo comunicado.

  \item[E2: alejamiento.] Misma geometría, pero con trayectoria divergente.
    No debe existir señal.

  \item[E3: cruce diagonal.] Aproximación diagonal a una esquina para evaluar
    simultáneamente distancia, velocidad y clasificación lateral.

  \item[E4: retroceso.] Obstáculo previamente observado detrás del usuario y
    desplazamiento hacia atrás.

  \item[E5: transición de FOV.] Giro de cabeza que hace que un obstáculo pase
    del lateral al campo frontal y viceversa.

  \item[E6: detención.] Frenado frente a un riesgo lateral. La intensidad debe
    caer al desaparecer la velocidad de cierre.
\end{description}

Las métricas recomendadas son latencia de activación, tasa de falsos
positivos, tasa de falsos negativos, continuidad de intensidad, error de TTC
frente a una referencia obtenida del rosbag y porcentaje de tiempo con datos
inválidos.

\subsection{Decisiones de diseño}
\label{subsec:sa-decisions}

\paragraph{Uso exclusivo de obstáculos fuera del campo comunicado.}
El módulo no duplica la señal frontal. El ángulo consumido representa el campo
que efectivamente llega al usuario, no necesariamente el campo físico de la
cámara.

\paragraph{Un mensaje con tres canales.}
Izquierda, derecha y atrás se publican juntos para mantener coherencia
temporal. Cada canal conserva intensidad y métricas diagnósticas, aunque el
actuador final solo requiera dirección e intensidad.

\paragraph{TTC como variable principal de riesgo.}
El TTC combina distancia y velocidad de cierre en una magnitud interpretable.
Evita ajustar dos factores independientes sin una justificación cinemática.

\paragraph{Velocidad de RTAB-Map con respaldo por pose.}
Se prioriza el campo \texttt{twist} por ser la interfaz convencional de
\texttt{nav\_msgs/Odometry}, pero no se asume que su frame o calidad sean
correctos sin verificación. La diferenciación de pose proporciona una ruta de
respaldo determinista.

\paragraph{Radio corporal como corredor, no como resta de distancia.}
Los límites de $0{,}20$ y $1{,}00\,\text{m}$ se aplican a la distancia
celda-cámara solicitada. El radio corporal determina si la trayectoria
intersecta espacialmente una celda y evita descontar dos veces la proximidad
del cuerpo.

\paragraph{Sin señal cuando el usuario está quieto.}
La proximidad por sí sola no activa el módulo. Su función es prevenir
colisiones causadas por el movimiento hacia obstáculos que ya no se están
comunicando.

\paragraph{Sin uso directo de velocidad angular.}
La orientación actual determina qué celdas quedan fuera del FOV. Una celda que
entra al campo frontal deja de señalizarse automáticamente; no es necesario
agregar la velocidad de giro a la intensidad.

\paragraph{No modificar el envejecimiento del mapa en esta etapa.}
La edad de observación por celda pertenece al modelo del
\texttt{local\_mapper}. \texttt{spatial\_awareness} consumirá el estado que el
mapa publique y solo verificará la antigüedad global del mensaje.

\subsection{Limitaciones y trabajo futuro}
\label{subsec:sa-future}

El modelo de TTC supone obstáculos estáticos. Una persona u objeto móvil
requeriría estimar también su velocidad relativa. Además, la odometría visual
puede degradarse en escenas sin textura, con iluminación adversa o movimiento
brusco, afectando tanto la posición del mapa como la velocidad estimada.

La ausencia de edad por celda puede mantener temporalmente obstáculos que ya
no existen. Esta limitación se resolverá posteriormente en
\texttt{local\_mapper}. También queda pendiente estudiar una representación de
la huella corporal no circular, el efecto de la posición real de la cámara
respecto del torso y curvas perceptuales de intensidad obtenidas con usuarios.

Finalmente, la frecuencia y latencia extremo a extremo deberán medirse en la
Raspberry Pi~5. El objetivo preliminar es evaluar el riesgo a 10~Hz, consumir
cada mapa generado por la cámara a 6~Hz y mantener una latencia de
procesamiento inferior a $100\,\text{ms}$ desde la publicación del mapa.

\subsection{Conclusión}
\label{subsec:sa-conclusion}

El módulo \texttt{spatial\_awareness} extiende la señalización frontal con una
advertencia cinemática para obstáculos conocidos fuera del campo comunicado.
Su diseño utiliza el mapa acumulado en \texttt{odom}, la pose y velocidad de
RTAB-Map, y la apertura efectiva de la grilla háptica. Solo se activa cuando
el usuario se aproxima a una agrupación ocupada dentro del intervalo de
$0{,}20$ a $1{,}00\,\text{m}$.

La salida separa izquierda, derecha y atrás, mientras que la intensidad deriva
del tiempo estimado hasta la colisión. La arquitectura mantiene desacopladas
la construcción del mapa, la evaluación de riesgo y la futura actuación
háptica, y deja explícitos los parámetros que deberán calibrarse mediante
escenas sintéticas y registros de caminata reales.
